% Continuous editable source/English for Volume I physical native27–30.
% Native27 receives the next paragraph of Article29; native30 ends within Article35.
% This continuation is outside the existing native2–26 partial PDF masters.
\SourceLeaf{27}

Ponamus itaque $a = \delta e$, $m = \delta f$, $t - u = \delta k$, eritque $e$ ad $f$
primus. Tum vero congruentiae propositae $\delta e x + \delta k \equiv 0 \ (\mathrm{mod.}\ \delta f)$
aequivalebit haec $ex + k \equiv 0 \ (\mathrm{mod.}\ f)$, i.\,e. quicunque ipsius $x$ valor huic
satisfaciat, etiam illi satisfaciet et vice versa. Manifesto enim $ex+k$ per $f$ dividi
poterit, quando $\delta e x + \delta k$ per $\delta f$ dividi potest, et vice versa. At
congruentiam $ex+k\equiv 0\ (\mathrm{mod.}\ f)$ supra solvere docuimus; unde simul patet, si $v$ sit
unus ex valoribus ipsius $x$, $x \equiv v \ (\mathrm{mod.}\ f)$ exhibere resolutionem completam
congruentiae propositae.

\Article{30}

Quando modulus est compositus, nonnumquam praestat sequenti methodo uti.

Sit modulus $=mn$, atque congruentia proposita $ax\equiv b$. Solvatur primo congruentia
haec secundum modulum $m$, ponamusque ei satisfieri, si $x\equiv v \ (\mathrm{mod.}\ \frac{m}{\delta})$,
designante $\delta$ divisorem communem maximum numerorum $m,a$. Iam manifestum est,
quemvis valorem ipsius $x$ congruentiae $ax\equiv b$ secundum modulum $mn$ satisfacientem
eidem etiam secundum modulum $m$ satisfacere debere: adeoque in forma $v+\frac{m}{\delta}x'$
contineri, designante $x'$ numerum indeterminatum, quamvis non vice versa omnes numeri in
forma $v+\frac{m}{\delta}x'$ contenti congruentiae secundum mod. $mn$ satisfaciant. Quomodo
autem $x'$ determinari debeat, ut $v+\frac{m}{\delta}x'$ fiat radix congruentiae
$ax\equiv b \ (\mathrm{mod.}\ mn)$, ex solutione congruentiae $\frac{am}{\delta}x'+av\equiv b\ (\mathrm{mod.}\ mn)$
deduci potest, cui aequivalet haec $\frac{a}{\delta}x'\equiv\frac{b-av}{m}\ (\mathrm{mod.}\ n)$. Hinc
colligitur solutionem congruentiae cuiuscunque primi gradus secundum modulum $mn$ reduci
posse ad solutionem duarum congruentiarum secundum modulum $m$ et $n$. Facile autem
perspicietur, si $n$ iterum sit productum e duobus factoribus, solutionem congruentiae
secundum modulum $n$ pendere a solutione duarum congruentiarum quarum moduli sint illi
factores. Generaliter solutio congruentiae secundum modulum compositum quemcumque pendet
a solutione aliarum congruentiarum, quarum moduli sunt factores illius numeri; hi autem,
si commodum esse videtur, ita semper accipi possunt, ut sint numeri primi.

\emph{Ex.} Si congruentia $19x\equiv 1 \ (\mathrm{mod.}\ 140)$ proponitur: solvatur primo secundum
modulum $2$, eritque $x\equiv 1\ (\mathrm{mod.}\ 2)$. Ponatur $x=1+2x'$, fietque $38x'\equiv -18\ (\mathrm{mod.}\ 140)$
cui aequivalet $19x'\equiv -9\ (\mathrm{mod.}\ 70)$. Si haec
\SourceLeaf{28}
iterum secundum modulum $2$ solvitur, fit $x'\equiv 1\ (\mathrm{mod.}\ 2)$ positoque
$x'=1+2x''$, fit $38x''\equiv -28\ (\mathrm{mod.}\ 70)$ sive $19x''\equiv -14\ (\mathrm{mod.}\ 35)$.
Haec secundum $5$ soluta dat $x''\equiv 4\ (\mathrm{mod.}\ 5)$, substitutoque $x''=4+5x'''$,
fit $95x'''\equiv -90\ (\mathrm{mod.}\ 35)$ sive $19x'''\equiv -18\ (\mathrm{mod.}\ 7)$.
Ex hac tandem sequitur, $x'''\equiv 2\ (\mathrm{mod.}\ 7)$, positoque $x'''=2+7x''''$
colligitur $x=59+140x''''$; quare $x\equiv 59\ (\mathrm{mod.}\ 140)$ est solutio completa
congruentiae propositae.

\Article{31}

Simili modo ut aequationis $ax=b$ radix per $\frac{b}{a}$ exprimitur, etiam congruentiae
$ax\equiv b$ radicem quamcunque per $\frac{b}{a}$ designabimus, congruentiae modulum,
distinctionis gratia, apponentes. Ita e.\,g. $\frac{19}{17}\ (\mathrm{mod.}\ 12)$ denotat quemvis
numerum, qui est $\equiv 11\ (\mathrm{mod.}\ 12)$\begingroup\renewcommand{\thefootnote}{*)}\footnote{id quod ex analogia per $\frac{11}{1}\ (\mathrm{mod.}\ 12)$ designari potest.}\endgroup. Generaliter ex praecedentibus patet,
$\frac{a}{b}\ (\mathrm{mod.}\ c)$ nihil reale significare (aut si quis malit aliquid imaginarii), si
$a,c$ habeant divisorem communem, qui ipsum $b$ non metiatur. At hoc casu excepto,
expressio $\frac{b}{a}\ (\mathrm{mod.}\ c)$ semper valores reales habebit, et quidem infinitos: hi vero
omnes secundum $c$ erunt congrui quando $a$ ad $c$ primus, aut secundum $\frac{c}{\delta}$,
quando $\delta$ numerorum $c,a$ divisor communis maximus.

Hae expressiones similem fere habent algorithmum ut fractiones vulgares. Aliquot
proprietates quae facile ex praecedentibus deduci possunt hic apponimus.

\hangindent=1.7em\hangafter=1
1.\quad Si secundum modulum $c$, $a\equiv\alpha$, $b\equiv\beta$ expressiones
$\frac{a}{b}\ (\mathrm{mod.}\ c)$ et $\frac{\alpha}{\beta}\ (\mathrm{mod.}\ c)$ sunt aequivalentes.

\hangindent=1.7em\hangafter=1
2.\quad $\frac{a\delta}{b\delta}\ (\mathrm{mod.}\ c\delta)$ et $\frac{a}{b}\ (\mathrm{mod.}\ c)$ sunt aequivalentes.

\hangindent=1.7em\hangafter=1
3.\quad $\frac{ak}{bk}\ (\mathrm{mod.}\ c)$ et $\frac{a}{b}\ (\mathrm{mod.}\ c)$ sunt aequivalentes quando $k$ ad $c$ est primus.

Multae aliae similes propositiones afferri possent: at quum nulli difficultati sint
obnoxiae, neque ad sequentia adeo necessariae, ad alia properamus.

\begin{center}\small\emph{De inveniendo numero secundum modulos datos residuis datis congruo.}\end{center}

\Article{32}

Problema quod magnum in sequentibus usum habebit, \emph{invenire omnes numeros, qui
secundum modulos quotcunque datos residua data praebent}, facile ex praecedentibus
solvi potest. Sint primo duo moduli $A,B$, secundum quos numerus
\SourceLeaf{29}
quaesitus, $z$, numeris $a,b$ respective congruus esse debeat. Omnes itaque valores
ipsius $z$ sub forma $Ax+a$ continentur, ubi $x$ est indeterminatus sed talis ut fiat
$Ax+a\equiv b\ (\mathrm{mod.}\ B)$. Quodsi iam numerorum $A,B$ divisor communis maximus est $\delta$,
resolutio completa huius congruentiae hanc habebit formam: $x\equiv v\ (\mathrm{mod.}\ \frac{B}{\delta})$
sive quod eodem redit, $x=v+\frac{kB}{\delta}$, denotante $k$ numerum integrum arbitrarium.
Hinc formula $Av+a+\frac{kAB}{\delta}$ omnes ipsius $z$ valores comprehendet, i.\,e.
$z\equiv Av+a\ (\mathrm{mod.}\ \frac{AB}{\delta})$ erit resolutio completa problematis. Si ad modulos
$A,B$ tertius accedit, $C$, secundum quem numerus quaesitus $z$ debet esse $\equiv c$,
manifesto eodem modo procedendum, quum binae priores conditiones in unicam iam sint
conflatae. Scilicet si numerorum $\frac{AB}{\delta}.$
$C$ divisor communis maximus $=\varepsilon$, atque congruentiae
$\frac{AB}{\delta}x+Av+a\equiv c\ (\mathrm{mod.}\ C)$ resolutio:
$x\equiv w\ (\mathrm{mod.}\ \frac{C}{\varepsilon})$, problema per congruentiam
$z\equiv\frac{ABw}{\delta}+Av+a\ (\mathrm{mod.}\ \frac{ABC}{\delta\varepsilon})$ complete erit
resolutum. Similiter procedendum, quotcunque moduli proponantur. Observari convenit
$\frac{AB}{\delta}$, $\frac{ABC}{\delta\varepsilon}$ esse numerorum $A,B$; et $A,B,C$
respective minimos communes dividuos, facileque inde perspicitur, quotcunque habeantur
moduli $A,B,C$ etc., si eorum minimus communis dividuus sit $M$, resolutionem completam
hanc formam habere, $z\equiv r\ (\mathrm{mod.}\ M)$. Ceterum quando ulla congruentiarum auxiliarium
est irresolubilis, problema impossibilitatem involvere concludendum est. Perspicuum
vero, hoc evenire non posse, quando omnes numeri $A,B,C$ etc. inter se sint primi.

\emph{Ex.} Sint numeri $A,B,C$; $a,b,c$; $504,35,16$; $17,-4,33$; hic duae conditiones ut
$z$ sit $\equiv 17\ (\mathrm{mod.}\ 504)$ et $\equiv -4\ (\mathrm{mod.}\ 35)$ unicae. ut sit
$\equiv 521\ \mathrm{mod.}\ 2520)$ aequivalent; ex qua cum hac:
$z\equiv 33\ (\mathrm{mod.}\ 16)$ coniuncta, promanat $z\equiv 3041\ (\mathrm{mod.}\ 5040)$.

\Article{33}

Quando omnes numeri $A,B,C$ etc. inter se sunt primi, constat. productum ex ipsis esse
minimum omnibus communem dividuum. In quo casu manifestum est, omnes congruentias
$z\equiv a\ (\mathrm{mod.}\ A)$; $z\equiv b\ (\mathrm{mod.}\ B)$ etc. unicae
$z\equiv r\ (\mathrm{mod.}\ R)$ prorsus aequivalere, denotante $R$ numerorum $A,B,C$ etc.
productum. Hinc vero vicissim sequitur, unicam conditionem $z\equiv r\ (\mathrm{mod.}\ R)$
in plures dissolvi posse; scilicet si $R$ quomodocunque in factores inter se primos
$A,B,C$ etc. resolvitur, conditiones $z\equiv r\ (\mathrm{mod.}\ A)$,
$z\equiv r\ (\mathrm{mod.}\ B)$, $z\equiv r\ (\mathrm{mod.}\ C)$, etc. propositum exhaurient.
Haec observatio methodum nobis aperit
\SourceLeaf{30}
non modo impossibilitatem, si quam forte conditiones propositae implicent, statim
detegendi, sed etiam calculum commodius atque concinnius instituendi.

\Article{34}

Sint ut supra conditiones propositae, ut sit $z\equiv a\ (\mathrm{mod.}\ A)$,
$z\equiv b\ (\mathrm{mod.}\ B)$, $z\equiv c\ (\mathrm{mod.}\ C)$. Resolvantur omnes moduli in
factores inter se primos, $A$ in $A'A''A'''$ etc.; $B$ in $B'B''B'''$ etc. etc. et quidem
ita ut numeri $A',A''$ etc. $B',B''$ etc. etc. sint aut primi, aut primorum potestates.
Si vero aliquis numerorum $A,B,C$ etc. iam per se est primus, aut primi potestas, nulla
resolutione in factores pro hocce opus est. Tum vero ex praecedentibus patescit, pro
conditionibus propositis hasce substitui posse: $z\equiv a\ (\mathrm{mod.}\ A')$,
$z\equiv a\ (\mathrm{mod.}\ A'')$, $z\equiv a\ (\mathrm{mod.}\ A''')$ etc.,
$z\equiv b\ (\mathrm{mod.}\ B')$, $z\equiv b\ (\mathrm{mod.}\ B'')$ etc. etc. Iam nisi omnes
numeri $A,B,C$ etc. fuerint inter se primi, ex.\,gr. si $A$ ad $B$ non primus, manifestum
est, omnes divisores primos ipsorum $A,B$ diversos esse non posse, sed inter factores
$A',A'',A'''$ etc. unum aut alterum esse debere, qui inter $B',B'',B'''$ etc. aut aequalem
aut multiplum aut submultiplum habeat. Si \emph{primo} $A'=B'$, conditiones
$z\equiv a\ (\mathrm{mod.}\ A')$, $z\equiv b\ (\mathrm{mod.}\ B')$ identicae esse debent,
sive $a\equiv b\ (\mathrm{mod.}\ A'\ \text{vel}\ B')$, quare alterutra reiici poterit.
Si vero non $a\equiv b\ (\mathrm{mod.}\ A')$, problema impossibilitatem implicat.
Si \emph{secundo} $B'$ multiplum ipsius $A'$, conditio $z\equiv a\ (\mathrm{mod.}\ A')$ in
hac $z\equiv b\ (\mathrm{mod.}\ B')$ contenta esse debet, sive haec
$z\equiv b\ (\mathrm{mod.}\ A')$ quae ex posteriori deducitur cum priori identica esse
debet. Unde sequitur conditionem $z\equiv a\ (\mathrm{mod.}\ A')$, nisi alteri repugnet
(in quo casu problema impossibile) reiici posse. Quando omnes conditiones superfluae
ita reiectae sunt, patet, omnes modulos ex his $A',A'',A'''$ etc., $B',B'',B'''$ etc. etc.
remanentes inter se primos fore: tum igitur de problematis possibilitate certi esse
et secundum praecepta ante data procedere possumus.

\Article{35}

\emph{Ex.} Si ut supra esse debet $z\equiv 17\ (\mathrm{mod.}\ 504)$,
$\equiv -4\ (\mathrm{mod.}\ 35)$, et $\equiv 33\ (\mathrm{mod.}\ 16)$; hae conditiones in
sequentes resolvi possunt, $z\equiv 17\ (\mathrm{mod.}\ 8)$, $\equiv 17\ (\mathrm{mod.}\ 9)$,
$\equiv 17\ (\mathrm{mod.}\ 7)$, $\equiv -4\ (\mathrm{mod.}\ 5)$,
$\equiv -4\ (\mathrm{mod.}\ 7)$, $\equiv 33\ (\mathrm{mod.}\ 16)$.
Ex his conditiones $z\equiv 17\ (\mathrm{mod.}\ 8)$, $z\equiv 17\ (\mathrm{mod.}\ 7)$ reiici
possunt, quum prior in conditione $z\equiv 33\ (\mathrm{mod.}\ 16)$ contineatur, posterior
vero cum hac $z\equiv -4\ (\mathrm{mod.}\ 7)$ sit identica; remanent itaque
% Native30 stops at remanent itaque; Article35 and Section II remain open.
% Native31 has not been read or inferred in this assignment.
